% Options for packages loaded elsewhere
\PassOptionsToPackage{unicode}{hyperref}
\PassOptionsToPackage{hyphens}{url}
\PassOptionsToPackage{dvipsnames,svgnames,x11names}{xcolor}
%
\documentclass[
  letterpaper,
  DIV=11,
  numbers=noendperiod]{scrreprt}

\usepackage{amsmath,amssymb}
\usepackage{iftex}
\ifPDFTeX
  \usepackage[T1]{fontenc}
  \usepackage[utf8]{inputenc}
  \usepackage{textcomp} % provide euro and other symbols
\else % if luatex or xetex
  \usepackage{unicode-math}
  \defaultfontfeatures{Scale=MatchLowercase}
  \defaultfontfeatures[\rmfamily]{Ligatures=TeX,Scale=1}
\fi
\usepackage{lmodern}
\ifPDFTeX\else  
    % xetex/luatex font selection
\fi
% Use upquote if available, for straight quotes in verbatim environments
\IfFileExists{upquote.sty}{\usepackage{upquote}}{}
\IfFileExists{microtype.sty}{% use microtype if available
  \usepackage[]{microtype}
  \UseMicrotypeSet[protrusion]{basicmath} % disable protrusion for tt fonts
}{}
\makeatletter
\@ifundefined{KOMAClassName}{% if non-KOMA class
  \IfFileExists{parskip.sty}{%
    \usepackage{parskip}
  }{% else
    \setlength{\parindent}{0pt}
    \setlength{\parskip}{6pt plus 2pt minus 1pt}}
}{% if KOMA class
  \KOMAoptions{parskip=half}}
\makeatother
\usepackage{xcolor}
\setlength{\emergencystretch}{3em} % prevent overfull lines
\setcounter{secnumdepth}{5}
% Make \paragraph and \subparagraph free-standing
\makeatletter
\ifx\paragraph\undefined\else
  \let\oldparagraph\paragraph
  \renewcommand{\paragraph}{
    \@ifstar
      \xxxParagraphStar
      \xxxParagraphNoStar
  }
  \newcommand{\xxxParagraphStar}[1]{\oldparagraph*{#1}\mbox{}}
  \newcommand{\xxxParagraphNoStar}[1]{\oldparagraph{#1}\mbox{}}
\fi
\ifx\subparagraph\undefined\else
  \let\oldsubparagraph\subparagraph
  \renewcommand{\subparagraph}{
    \@ifstar
      \xxxSubParagraphStar
      \xxxSubParagraphNoStar
  }
  \newcommand{\xxxSubParagraphStar}[1]{\oldsubparagraph*{#1}\mbox{}}
  \newcommand{\xxxSubParagraphNoStar}[1]{\oldsubparagraph{#1}\mbox{}}
\fi
\makeatother

\usepackage{color}
\usepackage{fancyvrb}
\newcommand{\VerbBar}{|}
\newcommand{\VERB}{\Verb[commandchars=\\\{\}]}
\DefineVerbatimEnvironment{Highlighting}{Verbatim}{commandchars=\\\{\}}
% Add ',fontsize=\small' for more characters per line
\usepackage{framed}
\definecolor{shadecolor}{RGB}{241,243,245}
\newenvironment{Shaded}{\begin{snugshade}}{\end{snugshade}}
\newcommand{\AlertTok}[1]{\textcolor[rgb]{0.68,0.00,0.00}{#1}}
\newcommand{\AnnotationTok}[1]{\textcolor[rgb]{0.37,0.37,0.37}{#1}}
\newcommand{\AttributeTok}[1]{\textcolor[rgb]{0.40,0.45,0.13}{#1}}
\newcommand{\BaseNTok}[1]{\textcolor[rgb]{0.68,0.00,0.00}{#1}}
\newcommand{\BuiltInTok}[1]{\textcolor[rgb]{0.00,0.23,0.31}{#1}}
\newcommand{\CharTok}[1]{\textcolor[rgb]{0.13,0.47,0.30}{#1}}
\newcommand{\CommentTok}[1]{\textcolor[rgb]{0.37,0.37,0.37}{#1}}
\newcommand{\CommentVarTok}[1]{\textcolor[rgb]{0.37,0.37,0.37}{\textit{#1}}}
\newcommand{\ConstantTok}[1]{\textcolor[rgb]{0.56,0.35,0.01}{#1}}
\newcommand{\ControlFlowTok}[1]{\textcolor[rgb]{0.00,0.23,0.31}{\textbf{#1}}}
\newcommand{\DataTypeTok}[1]{\textcolor[rgb]{0.68,0.00,0.00}{#1}}
\newcommand{\DecValTok}[1]{\textcolor[rgb]{0.68,0.00,0.00}{#1}}
\newcommand{\DocumentationTok}[1]{\textcolor[rgb]{0.37,0.37,0.37}{\textit{#1}}}
\newcommand{\ErrorTok}[1]{\textcolor[rgb]{0.68,0.00,0.00}{#1}}
\newcommand{\ExtensionTok}[1]{\textcolor[rgb]{0.00,0.23,0.31}{#1}}
\newcommand{\FloatTok}[1]{\textcolor[rgb]{0.68,0.00,0.00}{#1}}
\newcommand{\FunctionTok}[1]{\textcolor[rgb]{0.28,0.35,0.67}{#1}}
\newcommand{\ImportTok}[1]{\textcolor[rgb]{0.00,0.46,0.62}{#1}}
\newcommand{\InformationTok}[1]{\textcolor[rgb]{0.37,0.37,0.37}{#1}}
\newcommand{\KeywordTok}[1]{\textcolor[rgb]{0.00,0.23,0.31}{\textbf{#1}}}
\newcommand{\NormalTok}[1]{\textcolor[rgb]{0.00,0.23,0.31}{#1}}
\newcommand{\OperatorTok}[1]{\textcolor[rgb]{0.37,0.37,0.37}{#1}}
\newcommand{\OtherTok}[1]{\textcolor[rgb]{0.00,0.23,0.31}{#1}}
\newcommand{\PreprocessorTok}[1]{\textcolor[rgb]{0.68,0.00,0.00}{#1}}
\newcommand{\RegionMarkerTok}[1]{\textcolor[rgb]{0.00,0.23,0.31}{#1}}
\newcommand{\SpecialCharTok}[1]{\textcolor[rgb]{0.37,0.37,0.37}{#1}}
\newcommand{\SpecialStringTok}[1]{\textcolor[rgb]{0.13,0.47,0.30}{#1}}
\newcommand{\StringTok}[1]{\textcolor[rgb]{0.13,0.47,0.30}{#1}}
\newcommand{\VariableTok}[1]{\textcolor[rgb]{0.07,0.07,0.07}{#1}}
\newcommand{\VerbatimStringTok}[1]{\textcolor[rgb]{0.13,0.47,0.30}{#1}}
\newcommand{\WarningTok}[1]{\textcolor[rgb]{0.37,0.37,0.37}{\textit{#1}}}

\providecommand{\tightlist}{%
  \setlength{\itemsep}{0pt}\setlength{\parskip}{0pt}}\usepackage{longtable,booktabs,array}
\usepackage{calc} % for calculating minipage widths
% Correct order of tables after \paragraph or \subparagraph
\usepackage{etoolbox}
\makeatletter
\patchcmd\longtable{\par}{\if@noskipsec\mbox{}\fi\par}{}{}
\makeatother
% Allow footnotes in longtable head/foot
\IfFileExists{footnotehyper.sty}{\usepackage{footnotehyper}}{\usepackage{footnote}}
\makesavenoteenv{longtable}
\usepackage{graphicx}
\makeatletter
\newsavebox\pandoc@box
\newcommand*\pandocbounded[1]{% scales image to fit in text height/width
  \sbox\pandoc@box{#1}%
  \Gscale@div\@tempa{\textheight}{\dimexpr\ht\pandoc@box+\dp\pandoc@box\relax}%
  \Gscale@div\@tempb{\linewidth}{\wd\pandoc@box}%
  \ifdim\@tempb\p@<\@tempa\p@\let\@tempa\@tempb\fi% select the smaller of both
  \ifdim\@tempa\p@<\p@\scalebox{\@tempa}{\usebox\pandoc@box}%
  \else\usebox{\pandoc@box}%
  \fi%
}
% Set default figure placement to htbp
\def\fps@figure{htbp}
\makeatother
% definitions for citeproc citations
\NewDocumentCommand\citeproctext{}{}
\NewDocumentCommand\citeproc{mm}{%
  \begingroup\def\citeproctext{#2}\cite{#1}\endgroup}
\makeatletter
 % allow citations to break across lines
 \let\@cite@ofmt\@firstofone
 % avoid brackets around text for \cite:
 \def\@biblabel#1{}
 \def\@cite#1#2{{#1\if@tempswa , #2\fi}}
\makeatother
\newlength{\cslhangindent}
\setlength{\cslhangindent}{1.5em}
\newlength{\csllabelwidth}
\setlength{\csllabelwidth}{3em}
\newenvironment{CSLReferences}[2] % #1 hanging-indent, #2 entry-spacing
 {\begin{list}{}{%
  \setlength{\itemindent}{0pt}
  \setlength{\leftmargin}{0pt}
  \setlength{\parsep}{0pt}
  % turn on hanging indent if param 1 is 1
  \ifodd #1
   \setlength{\leftmargin}{\cslhangindent}
   \setlength{\itemindent}{-1\cslhangindent}
  \fi
  % set entry spacing
  \setlength{\itemsep}{#2\baselineskip}}}
 {\end{list}}
\usepackage{calc}
\newcommand{\CSLBlock}[1]{\hfill\break\parbox[t]{\linewidth}{\strut\ignorespaces#1\strut}}
\newcommand{\CSLLeftMargin}[1]{\parbox[t]{\csllabelwidth}{\strut#1\strut}}
\newcommand{\CSLRightInline}[1]{\parbox[t]{\linewidth - \csllabelwidth}{\strut#1\strut}}
\newcommand{\CSLIndent}[1]{\hspace{\cslhangindent}#1}

\usepackage{booktabs}
\usepackage{longtable}
\usepackage{array}
\usepackage{multirow}
\usepackage{wrapfig}
\usepackage{float}
\usepackage{colortbl}
\usepackage{pdflscape}
\usepackage{tabu}
\usepackage{threeparttable}
\usepackage{threeparttablex}
\usepackage[normalem]{ulem}
\usepackage{makecell}
\usepackage{xcolor}
\KOMAoption{captions}{tableheading}
\makeatletter
\@ifpackageloaded{bookmark}{}{\usepackage{bookmark}}
\makeatother
\makeatletter
\@ifpackageloaded{caption}{}{\usepackage{caption}}
\AtBeginDocument{%
\ifdefined\contentsname
  \renewcommand*\contentsname{Table of contents}
\else
  \newcommand\contentsname{Table of contents}
\fi
\ifdefined\listfigurename
  \renewcommand*\listfigurename{List of Figures}
\else
  \newcommand\listfigurename{List of Figures}
\fi
\ifdefined\listtablename
  \renewcommand*\listtablename{List of Tables}
\else
  \newcommand\listtablename{List of Tables}
\fi
\ifdefined\figurename
  \renewcommand*\figurename{Figure}
\else
  \newcommand\figurename{Figure}
\fi
\ifdefined\tablename
  \renewcommand*\tablename{Table}
\else
  \newcommand\tablename{Table}
\fi
}
\@ifpackageloaded{float}{}{\usepackage{float}}
\floatstyle{ruled}
\@ifundefined{c@chapter}{\newfloat{codelisting}{h}{lop}}{\newfloat{codelisting}{h}{lop}[chapter]}
\floatname{codelisting}{Listing}
\newcommand*\listoflistings{\listof{codelisting}{List of Listings}}
\makeatother
\makeatletter
\makeatother
\makeatletter
\@ifpackageloaded{caption}{}{\usepackage{caption}}
\@ifpackageloaded{subcaption}{}{\usepackage{subcaption}}
\makeatother

\usepackage{bookmark}

\IfFileExists{xurl.sty}{\usepackage{xurl}}{} % add URL line breaks if available
\urlstyle{same} % disable monospaced font for URLs
\hypersetup{
  pdftitle={CORN3000-EI},
  pdfauthor={Michael Hunt},
  colorlinks=true,
  linkcolor={blue},
  filecolor={Maroon},
  citecolor={Blue},
  urlcolor={Blue},
  pdfcreator={LaTeX via pandoc}}


\title{CORN3000-EI}
\author{Michael Hunt}
\date{Invalid Date}

\begin{document}
\maketitle

\renewcommand*\contentsname{Table of contents}
{
\hypersetup{linkcolor=}
\setcounter{tocdepth}{2}
\tableofcontents
}

\bookmarksetup{startatroot}

\chapter*{Preface}\label{preface}
\addcontentsline{toc}{chapter}{Preface}

\markboth{Preface}{Preface}

Here you will find course materials for the module CORN3000
Environmental Investigations. This is part of the BSc Environmental
Resource Management, delivered at CCUC: Newquay. This qualification is
validated by Plymouth University.

\bookmarksetup{startatroot}

\chapter{Gannel ERM 2025-2026
analysis}\label{gannel-erm-2025-2026-analysis}

Analysis of data from file trip to the Gannel 2/3/2026 with ERM group.

\section{Load packages}\label{load-packages}

\begin{Shaded}
\begin{Highlighting}[]
\FunctionTok{library}\NormalTok{(tidyverse)}
\FunctionTok{library}\NormalTok{(here)}
\FunctionTok{library}\NormalTok{(janitor) }\CommentTok{\# cleans names of columns}
\FunctionTok{library}\NormalTok{(readxl)}
\FunctionTok{library}\NormalTok{(cowplot)}
\end{Highlighting}
\end{Shaded}

\section{Load data}\label{load-data}

\begin{Shaded}
\begin{Highlighting}[]
\NormalTok{filepath }\OtherTok{\textless{}{-}} \FunctionTok{here}\NormalTok{(}\StringTok{"data"}\NormalTok{,}\StringTok{"gannel\_2526\_data.xlsx"}\NormalTok{)}
\NormalTok{gannel }\OtherTok{\textless{}{-}} \FunctionTok{read\_excel}\NormalTok{(filepath,}\AttributeTok{sheet =} \StringTok{"Sediment"}\NormalTok{)}
\CommentTok{\# glimpse(gannel)}
\end{Highlighting}
\end{Shaded}

\section{Organic matter per cubic metre of wet
sediment}\label{organic-matter-per-cubic-metre-of-wet-sediment}

We want to calculate the mass of organic matter per cubic metre of
sediment. This is given by \[
\text{kg OM/m}^3 = \rho_{ws} \times p_{ds} \times p_{om}
\] where \(\rho_{ws}\) is the density of the wet sediment, \(p_{ds}\) is
the proportion of dry sediment in the wet sediment and \(p_{om}\) is the
proportion of organic matter in the dry sediment.

We will calculate eeach of these terms separately

\section{Calculate proportions of water and dry sediment in the wet
sediment}\label{calculate-proportions-of-water-and-dry-sediment-in-the-wet-sediment}

\begin{Shaded}
\begin{Highlighting}[]
\CommentTok{\# proportion by mass of wet sediment that is dry sediment}
\NormalTok{gannel }\OtherTok{\textless{}{-}}\NormalTok{ gannel }\SpecialCharTok{|\textgreater{}}
  \FunctionTok{mutate}\NormalTok{ (}\AttributeTok{p\_ds =}\NormalTok{ (sediment\_cooked\_inc\_boat }\SpecialCharTok{{-}}\NormalTok{ boat\_empty )}\SpecialCharTok{/}\NormalTok{(boat\_wet\_sediment  }\SpecialCharTok{{-}}\NormalTok{ boat\_empty ))}
\end{Highlighting}
\end{Shaded}

\begin{Shaded}
\begin{Highlighting}[]
\CommentTok{\# proportion by mass of dry sediment that is organic matter}
\NormalTok{gannel }\OtherTok{\textless{}{-}}\NormalTok{ gannel }\SpecialCharTok{|\textgreater{}}
  \FunctionTok{mutate}\NormalTok{ (}\AttributeTok{p\_om =} \DecValTok{1}\SpecialCharTok{{-}}\NormalTok{  (crucible\_ash }\SpecialCharTok{{-}}\NormalTok{ crucible\_empty )}\SpecialCharTok{/}\NormalTok{(crucible\_sediment  }\SpecialCharTok{{-}}\NormalTok{ crucible\_empty ))}
\end{Highlighting}
\end{Shaded}

\section{Density of wet sediment}\label{density-of-wet-sediment}

To calculate wet sediment density, we note:

\[
\begin{align}
\frac{1}{\rho_{ws}} &= \frac{p_{ds}}{\rho_{ds}} + \frac{p_w}{\rho_w}\\
&=\frac{p_{ds}\rho_{w} + p_w\rho_{ds}}{\rho_{ds}\rho_w}
\end{align}
\] Hence

\[
\begin{align}
\rho_{ws}=\frac{\rho_{ds}\rho_w}{p_{ds}\rho_{w} + p_w\rho_{ds}}
\end{align}
\] But \(p_w = 1-p_{ds}\) and \(\rho_w=1000\)kg/m\textsuperscript{3}.

Hence \[
\begin{align}
\rho_{ws}=\frac{1000\rho_{ds}}{1000p_{ds} + (1-p_{ds})\rho_{ds}}
\end{align}
\] Check that this formula makes sense: what does it give as the density
of the wet density if the proportion of dry sediment is either zero or
one?

We see that we need the density of the dry sediment. We measured the
quantities need to calculate this for our samples:

\begin{Shaded}
\begin{Highlighting}[]
\CommentTok{\# density of dry sediment}
\NormalTok{gannel }\OtherTok{\textless{}{-}}\NormalTok{ gannel }\SpecialCharTok{|\textgreater{}}
\FunctionTok{mutate}\NormalTok{(}\AttributeTok{rho\_ds =} \DecValTok{1000}\SpecialCharTok{*}\NormalTok{ soil\_mass}\SpecialCharTok{/}\NormalTok{water\_displacement)}
\end{Highlighting}
\end{Shaded}

We can now calcuate the density of wet sediment, given the density of
dry sediment \(\rho_{ds}\) and the proportion \(p_{ds}\) of wet sediment
that is dry sediment:

\begin{Shaded}
\begin{Highlighting}[]
\NormalTok{gannel }\OtherTok{\textless{}{-}}\NormalTok{ gannel }\SpecialCharTok{|\textgreater{}}
  \FunctionTok{mutate}\NormalTok{(}\AttributeTok{rho\_ws =} \DecValTok{1000}\SpecialCharTok{*}\NormalTok{rho\_ds }\SpecialCharTok{/}\NormalTok{ (}\DecValTok{1000}\SpecialCharTok{*}\NormalTok{p\_ds }\SpecialCharTok{+}\NormalTok{ (}\DecValTok{1}\SpecialCharTok{{-}}\NormalTok{p\_ds)}\SpecialCharTok{*}\NormalTok{rho\_ds))}
\end{Highlighting}
\end{Shaded}

Now that we have all the terms we require, we can calculate the mass of
organic matter per cubic metre of wet sediment:

\begin{Shaded}
\begin{Highlighting}[]
\NormalTok{gannel }\OtherTok{\textless{}{-}}\NormalTok{ gannel }\SpecialCharTok{|\textgreater{}}
  \FunctionTok{mutate}\NormalTok{(}\AttributeTok{mass\_om =}\NormalTok{ rho\_ws }\SpecialCharTok{*}\NormalTok{ p\_ds }\SpecialCharTok{*}\NormalTok{ p\_om)}
\end{Highlighting}
\end{Shaded}

The carbon content is about 0.58 times the organic matter content
(citation needed), and the mass of carbon per hectare in tonnes is
numerically 10 times the mass of carnon per cubic metre in kg. Hence

\begin{Shaded}
\begin{Highlighting}[]
\NormalTok{gannel }\OtherTok{\textless{}{-}}\NormalTok{ gannel }\SpecialCharTok{|\textgreater{}}
  \FunctionTok{mutate}\NormalTok{(}\AttributeTok{mass\_c =}\FloatTok{0.58} \SpecialCharTok{*} \DecValTok{10} \SpecialCharTok{*}\NormalTok{ mass\_om)}
\end{Highlighting}
\end{Shaded}

We can now find the means of each of the two values measured for each
transect

\begin{Shaded}
\begin{Highlighting}[]
\NormalTok{carbon\_summary }\OtherTok{\textless{}{-}}\NormalTok{ gannel}\SpecialCharTok{|\textgreater{}}
  \FunctionTok{group\_by}\NormalTok{(transect) }\SpecialCharTok{|\textgreater{}}
  \FunctionTok{summarise}\NormalTok{(}\AttributeTok{mean\_c\_mass =} \FunctionTok{mean}\NormalTok{(mass\_c))}
\end{Highlighting}
\end{Shaded}

Now we plot it:

\begin{Shaded}
\begin{Highlighting}[]
\NormalTok{carbon\_summary }\SpecialCharTok{|\textgreater{}}
  \FunctionTok{filter}\NormalTok{(transect}\SpecialCharTok{\textless{}}\DecValTok{8}\NormalTok{) }\SpecialCharTok{|\textgreater{}}
  \FunctionTok{ggplot}\NormalTok{(}\FunctionTok{aes}\NormalTok{(}\AttributeTok{x =}\NormalTok{ transect, }\AttributeTok{y =}\NormalTok{ mean\_c\_mass)) }\SpecialCharTok{+}
  \FunctionTok{geom\_point}\NormalTok{() }\SpecialCharTok{+}
  \FunctionTok{geom\_smooth}\NormalTok{(}\AttributeTok{method =} \StringTok{"lm"}\NormalTok{) }\SpecialCharTok{+}
  \FunctionTok{labs}\NormalTok{(}\AttributeTok{x =} \StringTok{"Transect"}\NormalTok{, }
       \AttributeTok{y =} \StringTok{"Soil carbon tonnes / hectare"}\NormalTok{) }\SpecialCharTok{+}
  \FunctionTok{scale\_y\_continuous}\NormalTok{(}\AttributeTok{limits =} \FunctionTok{c}\NormalTok{(}\DecValTok{0}\NormalTok{,}\DecValTok{700}\NormalTok{)) }\SpecialCharTok{+}
  \FunctionTok{theme\_cowplot}\NormalTok{()}
\end{Highlighting}
\end{Shaded}

\includegraphics[width=0.8\linewidth,height=\textheight,keepaspectratio]{gannel_ERM_2526_files/figure-pdf/unnamed-chunk-11-1.pdf}

A linear regression:

\begin{Shaded}
\begin{Highlighting}[]
\NormalTok{c\_model }\OtherTok{\textless{}{-}}\FunctionTok{lm}\NormalTok{(mean\_c\_mass }\SpecialCharTok{\textasciitilde{}}\NormalTok{ transect, }\AttributeTok{data =}\NormalTok{ carbon\_summary)}
\FunctionTok{summary}\NormalTok{(c\_model)}
\end{Highlighting}
\end{Shaded}

\begin{verbatim}

Call:
lm(formula = mean_c_mass ~ transect, data = carbon_summary)

Residuals:
   Min     1Q Median     3Q    Max 
-125.7  -73.5  -20.9   59.7  183.7 

Coefficients:
            Estimate Std. Error t value Pr(>|t|)  
(Intercept)    282.0       87.1    3.24    0.018 *
transect        36.9       17.2    2.14    0.076 .
---
Signif. codes:  0 '***' 0.001 '**' 0.01 '*' 0.05 '.' 0.1 ' ' 1

Residual standard error: 112 on 6 degrees of freedom
Multiple R-squared:  0.432, Adjusted R-squared:  0.338 
F-statistic: 4.57 on 1 and 6 DF,  p-value: 0.0765
\end{verbatim}

\begin{itemize}
\tightlist
\item
  It would be better to use distance to the shore line as the
  explanatory variable, instead of transect number
\item
  Is there evidence from this regression analysis of a change in carbon
  content from one transect to another?
\item
  What is the mean carbon content stored per hectare in the top metre
  depth of sediment, according to these data?
\item
  How does this compare with values in the literature, for salt marshes?
\item
  How does it compare with the soils of other habitats such as
  woodlands?
\end{itemize}

\bookmarksetup{startatroot}

\chapter*{References}\label{references}
\addcontentsline{toc}{chapter}{References}

\markboth{References}{References}

\phantomsection\label{refs}
\begin{CSLReferences}{0}{1}
\end{CSLReferences}




\end{document}
